\section{Methods}

\subsection{Interactome and modules}
Human protein functional associations were obtained from the STRING v12.0 functional association network \citep{Szklarczyk2023} (combined score integrating experimental, curated-database, co-expression, genomic-context, and text-mining channels; we use ``interactome'' in this broad association-network sense, not as a purely physical binding network) at combined-confidence $\geq$900, mapped to HUGO gene symbols, and filtered to genes with GTEx v8 liver median TPM $\geq$1 \citep{GTEx2020}; the largest connected component (LCC) was retained using NetworkX \citep{Hagberg2008} (7{,}677 nodes, 66{,}908 edges, from a raw $\geq$900 STRING network of 11{,}693 nodes / 100{,}383 edges; the $\geq$700 network, 9{,}773 LCC nodes, is used for threshold robustness). DILI-associated genes were obtained from DisGeNET curated associations (UMLS C0860207) \citep{Pinero2020}, giving 82 genes in the $\geq$900 LCC; this is a curated association module, not a validated causal gene set. Hyperforin targets (10 in the LCC) were curated from primary literature \citep{Moore2000, Watkins2001, Obach2000, Komoroski2004, Hennessy2002, Assefa2004, Wang2004, Quiney2007, Quiney2006}; Quercetin targets (62 in the LCC) were assembled from ChEMBL-derived bioactivity records \citep{Mendez2019} (the programmatic retrieval is documented in the repository's data manifest rather than re-executed here; see Data Availability). Protein identifiers were mapped to HUGO symbols via UniProt \citep{UniProt2023}, with non-human proteins excluded and symbols standardized (e.g.\ MDR1$\rightarrow$ABCB1); raw counts (14 Hyperforin, 122 Quercetin) reduce to 10 and 62 in the LCC (Table~\ref{tab:s_curation}). Five genes are targeted by both compounds (ABCG2, AKT1, CYP3A4, MMP2, MMP9) and are retained in both sets; four further literature-reported Hyperforin targets (TRPC6, VEGFA, PMAIP1, GRIN1) fall outside the liver-expressed LCC and are not analyzed \citep{Leuner2007, Hostanska2003, Kumar2006}.

\subsection{Closest distance and proximity Z-score}
Closest distance is the mean shortest path from each target to the nearest disease-module gene \citep{Guney2016}:
\begin{equation}
d_c(T,D) = \frac{1}{|T|}\sum_{t\in T}\min_{s\in D}\mathrm{dist}(t,s).
\label{eq:dc}
\end{equation}
The proximity Z-score standardizes this raw distance against a degree-matched null:
\begin{equation}
Z = \frac{d_c - \mu_{\text{null}}}{\sigma_{\text{null}}},
\label{eq:proxz}
\end{equation}
with $Z=0$ when $\sigma_{\text{null}}=0$. We use $n=1{,}000$ degree-matched random reference sets, fixed seed 42. Following Guney et al.\ \citep{Guney2016} we report $Z$ as the primary magnitude; we additionally report empirical, directional permutation p-values using the conservative $(r+1)/(n+1)$ convention \citep{Phipson2010} --- one-sided for smaller-than-null distance (proximity) and larger-than-null influence (RWR/EWI) --- whose floor at $n=1{,}000$ is $9.99\times10^{-4}$ (reported as $p<0.001$). For both compounds, $p<0.001$ confirms non-random proximity to the DILI module but does not discriminate between them; the Z-score is the primary comparative statistic. No two-tailed permutation p-value is used. Degree matching is performed both with a $\pm$25\% degree window (primary) and with Guney's $\geq$100-node degree-binning; the degree-binning choice has minimal effect under the fixed-disease null (Table~\ref{tab:s_guney}). In the Guney-fidelity comparison, we use random seed 42 rather than the toolbox default (452456), and the local shortest-path comparison sampler permits the original seed node in its own candidate pool whereas the toolbox excludes it; this is a $\leq$1\% candidate-pool difference for bins of $\geq$100 nodes, and the resulting Z-scores are unchanged within permutation tolerance (Table~\ref{tab:s_guney}). The shared RWR/EWI permutation sampler excludes original targets from the random candidate pool. The fixed-disease null (randomizing targets only) is primary, as it corresponds to comparing two compounds against one disease module; the two-sided null (randomizing both targets and disease, as in the canonical implementation) is reported as a robustness check.

\subsection{Random walk with restart and perturbation efficiency}
RWR uses $W=AD^{-1}$, restart $\alpha=0.15$ (a PageRank-style damping value; the canonical RWR reference is K\"ohler et al.\ \citep{Kohler2008}, where $\alpha$ is typically larger and results are reported robust to it; Guney et al.\ \citep{Guney2016} use shortest-path proximity and define no restart parameter), convergence tolerance $10^{-6}$, and $\leq$100 iterations. Influence and perturbation efficiency are defined by Eqs.~\eqref{eq:rwr}--\eqref{eq:pe}; perturbation efficiency equals mean per-target influence by linearity of the steady state in the restart vector. Degree-matched permutation nulls ($n=1{,}000$, seed 42) provide influence Z-scores and empirical p-values; Benjamini--Hochberg correction is applied within each network threshold. A full restart-probability sweep ($\alpha\in[0.10,0.70]$) is reported.

\subsection{Expression weighting}
Transitions are weighted by destination-node liver expression, $W'_{ij} = A_{ij} e_i / \sum_k A_{kj} e_k$, with $e_i$ the min--max-normalized $\log_2(\mathrm{TPM}+1)$ from GTEx v8 liver \citep{GTEx2020}; an expression floor of 0.01 maintains reachability and is shown to be immaterial (floor sweep $0$--$0.1$).

\subsection{Direct--propagated decomposition}
To separate direct target--disease overlap from propagated influence we isolate the propagated component by leave-one-out (scoring each target against the module excluding the target itself, thereby excluding the disease-member seed's own contribution from its module score) and report the following controls: removal of DILI-member seeds; removal of cytochrome-P450/transporter seeds; removal of targets shared by both compounds; comparison to global degree-matched random 10-gene sets; and size-matched Quercetin subsets. The network separation measure $S_{AB}$ \citep{Menche2015} is reported as an orthogonal overlap statistic. As a network-construction robustness check, we additionally rebuilt the interactome excluding STRING edges supported solely by text mining \citep{Huang2018} and re-ran all metrics, arguing against text-mining-only edges as the driver of the observed effect ordering (Supplementary Note, \textit{Text-mining robustness}).

\subsection{Chemical similarity}
Morgan fingerprints (ECFP4; radius 2, 2048 bits) were generated with RDKit \citep{RDKit2023}; maximum Tanimoto similarity to DILIrank 2.0 reference drugs \citep{Chen2016, Olubamiwa2025} was computed, with 0.4 as the structural-analogue threshold \citep{Maggiora2014}. DILIrank 2.0 classifies 1{,}336 drugs into four concern categories; we used the binary subset with retrievable PubChem SMILES, mapping vMost- and vLess-DILI-concern to DILI-positive (542 drugs) and vNo-DILI-concern to DILI-negative (365 drugs), and excluding the ambiguous-DILI-concern category.

\subsection{Operating-regime benchmark}
To assess whether the statistical mechanism underlying the effect-size/evidence dissociation extends beyond the worked pair, we sampled random target-set probes on a ten-point size grid ($|T|=5$--$80$, including the observed 10- and 62-target sizes; 20{,}000 probes per size) from the liver LCC and computed each probe's closest distance to the DILI module and its size-level calibration Z-score. Primary probes were degree-distribution-pinned to the full compound-target row degree profile (10 Hyperforin rows plus 62 Quercetin rows) using Guney's $\geq$100-node bins: for each probe, degree-bin labels were sampled with replacement from this target-row degree-bin distribution, and nodes were then sampled from the corresponding bins after excluding DILI-module genes from the candidate pool so distance-zero mass cannot scale with $|T|$. A non-DILI-target-profile family and a uniform-from-pool family were used as contrasts. The null parameters ($\mu_{\text{null}}$, $\sigma_{\text{null}}$) are estimated once per (size, degree-profile) family from the 20{,}000-probe reference distribution itself; each probe's calibration Z-score is its deviation from that shared size-level reference, not a separate per-probe permutation null. Because $N=20{,}000$, self-inclusion has negligible influence on size-level null parameters. We fitted the null-SD scaling exponent (log--log, bootstrap CI) and compared it across the DILI module and three random size-matched pseudo-modules. Rank discordance was quantified on cross-size probe pairs (500{,}000 per size pair; smaller set fixed at $|T|=10$): the maximum overturnable raw margin $\delta_{\max}(R)=(\mu_L-\mu_S)+|z_S|\,\sigma_L(\sqrt{R}-1)$, and the margin-conditional reversal rate (fraction of pairs in which the smaller set is closer by $\geq\delta_0$ hops yet the larger set has the more negative calibration Z-score) for pre-specified descriptive cutoffs $\delta_0\in\{0.3,0.5\}$, with a no-shrinkage ($\sigma_L:=\sigma_S$) counterfactual. These $\delta_0$ values summarize material-margin behavior in the probe benchmark; they are not used to select or validate the worked pair, which is located separately by its observed margin, margin percentile, and $\delta_{\max}$ envelope. Unconditional quantities are reported only as descriptive contrasts, with the H/Q-directional pattern (smaller set closer and larger set stronger by Z) distinguished from bidirectional rank discordance. This is a calibration benchmark of the statistic, not a drug-level toxicity predictor. Fixed seed 42; \texttt{scripts/run\_operating\_regime\_benchmark.py}, with results in \texttt{results/tables/operating\_regime\_*.csv}.

\subsection{Reproducibility}
All analyses run from the committed processed data with fixed seed 42, using Python 3.12 (NetworkX 3.6, NumPy 2.5, pandas 2.3, SciPy 1.18) and RDKit 2026.03; figures use R 4.6 with ggplot2 4.0. Exact pinned versions are in \texttt{requirements-lock.txt}. The committed pipeline and reviewer-evidence scripts regenerate the Results numbers from the deposited data.

\subsection*{Use of AI tools}
AI-assisted tools were used for code drafting and language editing. All analyses, results, and interpretations were verified by the author.
